\chapter{2010年大纲II卷（文）}

\section{单选题}

\begin{problem}
设全集 \(U = \{x \in \N^* \mid x < 6\}\)，集合 \(A = \{1, 3\}\)，\(B = \{3, 5\}\)，则 \(\complement_U\left( A \cup B \right) =\)（\quad）

\choices
    {\(\{1, 4\}\)}
    {\(\{1, 5\}\)}
    {\(\{2, 4\}\)}
    {\(\{2,5\}\)}

\begin{answer}
C
\end{answer}

\begin{solution}
由 \(x\in\N^*\) 且 \(x<6\)，得 \(U=\{1,2,3,4,5\}\). 又 \(A\cup B=\{1,3,5\}\)，所以 \(\complement_U\left(A\cup B\right)=\{2,4\}\). 故选 C.
\end{solution}
\end{problem}

\begin{problem}
不等式 \(\frac{x - 3}{x + 2} < 0\) 的解集为（\quad）

\choices
    {\(\{x \mid -2 < x < 3\}\)}
    {\(\{x \mid x < -2\}\)}
    {\(\{x \mid x < -2 \text{ 或 } x > 3\}\)}
    {\(\{x \mid x > 3\}\)}

\begin{answer}
A
\end{answer}

\begin{solution}
不等式有意义要求 \(x\ne-2\). 当 \(x<-2\) 时，\(x-3<0\)、\(x+2<0\)，分式为正；当 \(-2<x<3\) 时，\(x-3<0\)、\(x+2>0\)，分式为负；当 \(x>3\) 时，分式为正. 因而解集为 \((-2,3)\)，故选 A.
\end{solution}
\end{problem}

\begin{problem}
已知 \(\sin \alpha = \frac{2}{3}\)，则 \(\cos\left(\pi - 2\alpha\right) =\)（\quad）

\choices
    {\(-\frac{\sqrt{5}}{3}\)}
    {\(-\frac{1}{9}\)}
    {\(\frac{1}{9}\)}
    {\(\frac{\sqrt{5}}{3}\)}

\begin{answer}
B
\end{answer}

\begin{solution}
由诱导公式和二倍角公式，\(\cos\left(\pi-2\alpha\right)=-\cos 2\alpha=-(1-2\sin^2\alpha)=-(1-2\cdot\frac{4}{9})=-\frac{1}{9}\). 故选 B.
\end{solution}
\end{problem}

\begin{problem}
函数 \(y = 1 + \ln\left(x - 1\right)\)（\(x > 1\)）的反函数是（\quad）

\choices
    {\(y = \e^{x + 1} - 1\)（\(x > 0\)）}
    {\(y = \e^{x - 1} + 1\)（\(x > 0\)）}
    {\(y = \e^{x + 1} - 1\)（\(x \in \R\)）}
    {\(y = \e^{x - 1} + 1\)（\(x \in \R\)）}

\begin{answer}
D
\end{answer}

\begin{solution}
由 \(y=1+\ln\left(x-1\right)\) 得 \(y-1=\ln\left(x-1\right)\)，从而 \(x=\e^{y-1}+1\). 原函数的值域为 \(\R\)，交换 \(x\)，\(y\) 后反函数为 \(y=\e^{x-1}+1\)，且反函数的定义域为 \(x\in\R\). 故选 D.
\end{solution}
\end{problem}

\begin{problem}
若变量 \(x\)，\(y\) 满足约束条件 \(\begin{cases}
x \geqslant -1,\\
y \geqslant x,\\
3x + 2y \leqslant 5,
\end{cases}\) 则 \(z = 2x + y\) 的最大值为（\quad）

\choices
    {\(1\)}
    {\(2\)}
    {\(3\)}
    {\(4\)}

\begin{answer}
C
\end{answer}

\begin{solution}
由约束条件，\(x\leqslant y\leqslant\frac{5-3x}{2}\)，并且可行域要求 \(-1\leqslant x\leqslant1\). 对固定的 \(x\)，目标函数 \(z=2x+y\) 随 \(y\) 增大而增大，因此 \(z\leqslant2x+\frac{5-3x}{2}=\frac{x+5}{2}\leqslant3\). 当 \(x=1\)，\(y=1\) 时三个约束均满足，且 \(z=3\)，所以最大值为 \(3\). 故选 C.
\end{solution}
\end{problem}

\begin{problem}
如果等差数列 \(\{a_{n}\}\) 中，\(a_{3} + a_{4} + a_{5} = 12\)，那么 \(a_{1} + a_{2} + \cdots + a_{7} =\)（\quad）

\choices
    {\(14\)}
    {\(21\)}
    {\(28\)}
    {\(35\)}

\begin{answer}
C
\end{answer}

\begin{solution}
等差数列中 \(a_3+a_5=2a_4\)，所以 \(a_3+a_4+a_5=3a_4=12\)，得 \(a_4=4\). 又 \(a_1+a_7=a_2+a_6=a_3+a_5=2a_4\)，故 \(a_1+a_2+\cdots+a_7=7a_4=28\). 故选 C.
\end{solution}
\end{problem}

\begin{problem}
若曲线 \(y = x^{2} + ax + b\) 在点 \((0, b)\) 处的切线方程是 \(x - y + 1 = 0\)，则（\quad）

\choices
    {\(a = 1\)，\(b = 1\)}
    {\(a = -1\)，\(b = 1\)}
    {\(a = 1\)，\(b = -1\)}
    {\(a = -1\)，\(b = -1\)}

\begin{answer}
A
\end{answer}

\begin{solution}
切点为 \((0,b)\)，代入直线 \(x-y+1=0\) 得 \(b=1\). 曲线的导数为 \(y'=2x+a\)，所以在 \(x=0\) 处的切线斜率为 \(a\). 直线 \(x-y+1=0\) 即 \(y=x+1\)，斜率为 \(1\)，从而 \(a=1\). 故选 A.
\end{solution}
\end{problem}

\begin{problem}
在三棱锥 \(S - ABC\) 中，底面 \(ABC\) 为边长等于 \(2\) 的等边三角形，\(SA\) 垂直于底面 \(ABC\)，\(SA = 3\)，那么直线 \(AB\) 与平面 \(SBC\) 所成角的正弦值为（\quad）

\choices
    {\(\frac{\sqrt{3}}{4}\)}
    {\(\frac{\sqrt{5}}{4}\)}
    {\(\frac{\sqrt{7}}{4}\)}
    {\(\frac{3}{4}\)}

\begin{answer}
D
\end{answer}

\begin{solution}
取 \(A=(0,0,0)\)，\(B=(2,0,0)\)，\(C=(1,\sqrt{3},0)\)，\(S=(0,0,3)\). 则平面 \(SBC\) 的两个方向向量可取 \(\overrightarrow{SB}=(2,0,-3)\) 和 \(\overrightarrow{SC}=(1,\sqrt{3},-3)\)，其法向量可取 \(\bs n=\overrightarrow{SB}\times\overrightarrow{SC}=(3\sqrt{3},3,2\sqrt{3})\). 直线 \(AB\) 的方向向量为 \(\overrightarrow{AB}=(2,0,0)\). 若直线 \(AB\) 与平面 \(SBC\) 所成的角为 \(\theta\)，则 \(\sin\theta=\frac{|\overrightarrow{AB}\cdot\bs n|}{|\overrightarrow{AB}|\,|\bs n|}=\frac{6\sqrt{3}}{2\cdot4\sqrt{3}}=\frac{3}{4}\). 故选 D.
\end{solution}
\end{problem}

\begin{problem}
将标号为 \(1\)，\(2\)，\(3\)，\(4\)，\(5\)，\(6\) 的 \(6\) 张卡片放入 \(3\) 个不同的信封中，若每个信封放 \(2\) 张，其中标号为 \(1\)，\(2\) 的卡片放入同一信封，则不同的放法共有（\quad）

\choices
    {\(12\text{ 种}\)}
    {\(18\text{ 种}\)}
    {\(36\text{ 种}\)}
    {\(54\text{ 种}\)}

\begin{answer}
B
\end{answer}

\begin{solution}
先从 \(3\) 个不同信封中选出一个放置卡片 \(1,2\)，有 \(3\) 种选法. 剩下的 \(4\) 张卡片放入剩下的两个有标号信封中，每个放 \(2\) 张；选定其中一个信封所放的 \(2\) 张即可，故有 \(\mathrm C_4^2=6\) 种. 因此总放法数为 \(3\mathrm C_4^2=18\)，故选 B.
\end{solution}
\end{problem}

\begin{problem}
\(\triangle ABC\) 中，点 \(D\) 在边 \(AB\) 上，\(CD\) 平分 \(\angle ACB\). 若 \(\overrightarrow{CB} = \bs{a}\)，\(\overrightarrow{CA} = \bs{b}\)，\(|\bs{a}| = 1\)，\(|\bs{b}| = 2\)，则 \(\overrightarrow{CD} =\)（\quad）

\choices
    {\(\frac{1}{3} \bs{a} + \frac{2}{3} \bs{b}\)}
    {\(\frac{2}{3} \bs{a} + \frac{1}{3} \bs{b}\)}
    {\(\frac{3}{5} \bs{a} + \frac{4}{5} \bs{b}\)}
    {\(\frac{4}{5} \bs{a} + \frac{3}{5} \bs{b}\)}

\begin{answer}
B
\end{answer}

\begin{solution}
由角平分线定理，\(\frac{AD}{DB}=\frac{AC}{CB}=\frac{|\bs b|}{|\bs a|}=2\)，所以 \(\overrightarrow{AD}=\frac{2}{3}\overrightarrow{AB}\). 又 \(\overrightarrow{AB}=\overrightarrow{CB}-\overrightarrow{CA}=\bs a-\bs b\)，于是 \(\overrightarrow{CD}=\overrightarrow{CA}+\overrightarrow{AD}=\bs b+\frac{2}{3}(\bs a-\bs b)=\frac{2}{3}\bs a+\frac{1}{3}\bs b\). 故选 B.
\end{solution}
\end{problem}

\begin{problem}
与正方体 \(ABCD - A_{1}B_{1}C_{1}D_{1}\) 的三条棱 \(AB\)，\(CC_{1}\)，\(A_{1}D_{1}\) 所在直线的距离相等的点（\quad）

\choices
    {有且只有 \(1\text{ 个}\)}
    {有且只有 \(2\text{ 个}\)}
    {有且只有 \(3\text{ 个}\)}
    {有无数个}

\begin{answer}
D
\end{answer}

\begin{solution}
不妨设正方体棱长为 \(1\)，建立空间直角坐标系：\(A=(0,0,0)\)，\(B=(1,0,0)\)，\(C=(1,1,0)\)，\(D=(0,1,0)\)，\(A_1=(0,0,1)\)，\(B_1=(1,0,1)\)，\(C_1=(1,1,1)\)，\(D_1=(0,1,1)\). 在线段 \(B_1D\) 上任取一点 \(P=(1-t,t,1-t)\)，其中 \(0\leqslant t\leqslant1\).

点 \(P\) 到直线 \(AB\) 的距离为 \(\sqrt{t^2+(1-t)^2}\). 点 \(P\) 到直线 \(CC_1\) 的距离为 \(\sqrt{t^2+(1-t)^2}\). 点 \(P\) 到直线 \(A_1D_1\) 的距离也为 \(\sqrt{t^2+(1-t)^2}\). 因而线段 \(B_1D\) 上每一点都满足题设条件，满足条件的点有无数个. 故选 D.
\end{solution}
\end{problem}

\begin{problem}
已知椭圆 \(C: \frac{x^2}{a^2} + \frac{y^2}{b^2} = 1\)（\(a > b > 0\)）的离心率为 \(\frac{\sqrt{3}}{2}\)，过右焦点 \(F\) 且斜率为 \(k\)（\(k > 0\)）的直线与 \(C\) 相交于 \(A\)，\(B\) 两点. 若 \(\overrightarrow{AF} = 3\overrightarrow{FB}\)，则 \(k =\)（\quad）

\choices
    {\(1\)}
    {\(\sqrt{2}\)}
    {\(\sqrt{3}\)}
    {\(2\)}

\begin{answer}
B
\end{answer}

\begin{solution}
设椭圆焦距为 \(2c\). 由离心率 \(\frac{c}{a}=\frac{\sqrt{3}}{2}\)，得 \(c=\frac{\sqrt{3}}{2}a\)，从而 \(b^2=a^2-c^2=\frac{a^2}{4}\)，椭圆方程可写为 \(x^2+4y^2=a^2\). 右焦点为 \(F=(c,0)\)，设直线与椭圆的交点 \(B\) 为 \(B=(c+t,kt)\). 由 \(\overrightarrow{AF}=3\overrightarrow{FB}\)，得 \(A=(c-3t,-3kt)\)，且 \(t\ne0\).

将 \(A\)，\(B\) 代入椭圆方程并相减，得 \((c-3t)^2+36k^2t^2-(c+t)^2-4k^2t^2=0\)，即 \(t=\frac{c}{1+4k^2}\). 再将 \(B\) 代入椭圆方程，得 \(c^2+2ct+(1+4k^2)t^2=a^2\). 令 \(q=1+4k^2\)，代入 \(c^2=\frac{3}{4}a^2\) 和 \(t=\frac{c}{q}\)，得到 \(\frac{3}{4}a^2\left(1+\frac{3}{q}\right)=a^2\)，所以 \(1+\frac{3}{q}=\frac{4}{3}\)，即 \(q=9\). 因而 \(1+4k^2=9\)，又 \(k>0\)，得 \(k=\sqrt{2}\). 故选 B.
\end{solution}
\end{problem}

\section{填空题}

\begin{problem}
已知 \(\alpha\) 是第二象限的角，\(\tan \alpha = -\frac{1}{2}\)，则 \(\cos \alpha =\fillinblank{}\).
\end{problem}

\begin{answer}
\(-\frac{2\sqrt{5}}{5}\)
\end{answer}

\begin{solution}
因 \(\alpha\) 是第二象限角，故 \(\cos\alpha<0\). 由 \(\tan\alpha=-\frac{1}{2}\)，可设 \(\sin\alpha=-\frac{1}{2}\cos\alpha\). 代入 \(\sin^2\alpha+\cos^2\alpha=1\)，得 \(\frac{1}{4}\cos^2\alpha+\cos^2\alpha=1\)，即 \(\cos^2\alpha=\frac{4}{5}\). 取第二象限对应的负值，得 \(\cos\alpha=-\frac{2\sqrt{5}}{5}\).
\end{solution}

\begin{problem}
\(\left(x + \frac{1}{x}\right)^9\) 的展开式中 \(x^3\) 的系数是\fillinblank{}.
\end{problem}

\begin{answer}
84
\end{answer}

\begin{solution}
展开式通项为 \(T_{r+1}=\mathrm C_9^r x^{9-r}\left(\frac{1}{x}\right)^r=\mathrm C_9^r x^{9-2r}\). 要得到 \(x^3\)，应有 \(9-2r=3\)，解得 \(r=3\). 因而 \(x^3\) 的系数为 \(\mathrm C_9^3=84\).
\end{solution}

\begin{problem}
已知抛物线 \(C: y^{2} = 2px\) （\(p > 0\)）的准线为 \(l\)，过 \(M(1,0)\) 且斜率为 \(\sqrt{3}\) 的直线与 \(l\) 相交于点 \(A\)，与 \(C\) 的一个交点为 \(B\). 若 \(\overrightarrow{AM} = \overrightarrow{MB}\)，则 \(p =\fillinblank{}\).
\end{problem}

\begin{answer}
2
\end{answer}

\begin{solution}
抛物线 \(y^2=2px\) 的准线为 \(x=-\frac{p}{2}\). 过 \(M=(1,0)\) 且斜率为 \(\sqrt{3}\) 的直线为 \(y=\sqrt{3}(x-1)\)，所以 \(A=\left(-\frac{p}{2},-\sqrt{3}\left(\frac{p}{2}+1\right)\right)\). 由 \(\overrightarrow{AM}=\overrightarrow{MB}\)，点 \(M\) 是 \(AB\) 的中点，故 \(B=2M-A=\left(2+\frac{p}{2},\sqrt{3}\left(\frac{p}{2}+1\right)\right)\).

把 \(B\) 代入抛物线方程，得 \(3\left(\frac{p}{2}+1\right)^2=2p\left(2+\frac{p}{2}\right)\)，即 \(3(p+2)^2=4p(p+4)\). 化简得 \(p^2+4p-12=0\)，所以 \(p=2\) 或 \(p=-6\). 由 \(p>0\)，得 \(p=2\).
\end{solution}

\begin{problem}
已知球 \(O\) 的半径为 \(4\)，圆 \(M\) 与圆 \(N\) 为该球的两个小圆，\(AB\) 为圆 \(M\) 与圆 \(N\) 的公共弦，\(AB = 4\). 若 \(OM = ON = 3\)，则两圆圆心的距离 \(MN =\fillinblank{}\).
\end{problem}

\begin{answer}
3
\end{answer}

\begin{solution}
设 \(P\) 为公共弦 \(AB\) 的中点. 因为球半径为 \(4\)，且 \(OM=ON=3\)，所以小圆 \(M\)，\(N\) 的半径均为 \(\sqrt{4^2-3^2}=\sqrt{7}\). 于是 \(MP=NP=\sqrt{7-2^2}=\sqrt{3}\). 又 \(OP\perp AB\)，所以 \(OP=\sqrt{4^2-2^2}=2\sqrt{3}\).

点 \(O\)，\(M\)，\(N\)，\(P\) 共面，且 \(OM\perp MP\). 在直角三角形 \(OMP\) 中，设从 \(M\) 到 \(OP\) 的高为 \(h\)，由面积相等得 \(\frac{1}{2}OM\cdot MP=\frac{1}{2}OP\cdot h\)，所以 \(h=\frac{3\cdot\sqrt{3}}{2\sqrt{3}}=\frac{3}{2}\). 由于 \(OM=ON\) 且 \(MP=NP\)，直线 \(OP\) 是 \(MN\) 的垂直平分线，故 \(MN=2h=3\).
\end{solution}

\section{解答题}

\begin{problem}
\(\triangle ABC\) 中，\(D\) 为边 \(BC\) 上的一点，\(BD = 33\)，\(\sin B = \frac{5}{13}\)，\(\cos \angle ADC = \frac{3}{5}\)，求 \(AD\).
\end{problem}

\begin{solution}
由 \(\cos\angle ADC=\frac{3}{5}>0\)，知 \(\angle ADC\) 为锐角，故 \(\sin\angle ADC=\frac{4}{5}\). 又因 \(D\) 在 \(BC\) 上，三角形 \(ABD\) 中 \(\angle BAD=\angle ADC-B\). 由 \(\sin B=\frac{5}{13}\) 且 \(B<\angle ADC<\frac{\pi}{2}\)，得 \(\cos B=\frac{12}{13}\). 因而 \(\sin\angle BAD=\sin\left(\angle ADC-B\right)=\frac{4}{5}\cdot\frac{12}{13}-\frac{3}{5}\cdot\frac{5}{13}=\frac{33}{65}\).

在三角形 \(ABD\) 中应用正弦定理，得 \(\frac{AD}{\sin B}=\frac{BD}{\sin\angle BAD}\)，所以 \(AD=33\cdot\frac{5}{13}\cdot\frac{65}{33}=25\).
\end{solution}

\begin{problem}
已知 \(\{a_{n}\}\) 是各项均为正数的等比数列，且 \(a_{1} + a_{2} = 2\left(\frac{1}{a_{1}} + \frac{1}{a_{2}}\right)\)，\(a_{3} + a_{4} + a_{5} = 64\left(\frac{1}{a_{3}} + \frac{1}{a_{4}} + \frac{1}{a_{5}}\right)\).

\begin{enumerate}
\item 求 \( \{a_{n}\} \) 的通项公式；
\item 设 \( b_n = \left(a_n + \frac{1}{a_n}\right)^2 \)，求数列 \( \{b_n\} \) 的前 \( n \) 项和 \( T_n \).
\end{enumerate}
\end{problem}

\begin{solution}
\begin{enumerate}
\item 设首项为 \(a_1=u>0\)，公比为 \(q>0\). 由第一组条件，\(u(1+q)=2\left(\frac{1}{u}+\frac{1}{uq}\right)=\frac{2(1+q)}{uq}\)，从而 \(u^2q=2\). 由第二组条件，\(u q^2(1+q+q^2)=64\cdot\frac{1+q+q^2}{u q^4}\)，因 \(1+q+q^2>0\)，得 \(u^2q^6=64\)，即 \(u q^3=8\). 联立 \(u^2q=2\) 与 \(u q^3=8\)，可得 \(q^5=32\). 由于 \(q>0\)，所以 \(q=2\)，再得 \(u=1\). 因而 \(a_n=2^{n-1}\).
\item 由 \(a_n=2^{n-1}\)，\(b_n=\left(2^{n-1}+2^{1-n}\right)^2=4^{n-1}+2+4^{1-n}\). 所以前 \(n\) 项和为 \(T_n=\sum_{k=1}^{n}4^{k-1}+2n+\sum_{k=1}^{n}4^{1-k}=\frac{4^n-1}{3}+2n+\frac{4}{3}\left(1-4^{-n}\right)\).
\end{enumerate}
\end{solution}

\begin{problem}
如图，直三棱柱 \(ABC - A_{1}B_{1}C_{1}\) 中，\(AC = BC\)，\(AA_{1} = AB\)，\(D\) 为 \(BB_{1}\) 的中点，\(E\) 为 \(AB_{1}\) 上的一点，\(AE = 3EB_{1}\).

\begin{enumerate}
\item 证明：\(DE\) 为异面直线 \(AB_{1}\) 与 \(CD\) 的公垂线.
\item 设异面直线 \(AB_{1}\) 与 \(CD\) 的夹角为 \(45^{\circ}\)，求二面角 \(A_{1}-AC_{1}-B_{1}\) 的大小.
\end{enumerate}

\bitmapfigure[max width=0.42\linewidth]{img/2010/syllabus_paper_2_liberal/q19_fig1.png}
\end{problem}

\begin{solution}
\begin{enumerate}
\item 设 \(AB=s>0\)，建立直角坐标系：\(A=(0,0,0)\)，\(B=(s,0,0)\)，\(C=\left(\frac{s}{2},h,0\right)\)，\(A_1=(0,0,s)\)，\(B_1=(s,0,s)\)，\(C_1=\left(\frac{s}{2},h,s\right)\). 由 \(D\) 为 \(BB_1\) 的中点及 \(AE=3EB_1\)，得 \(D=\left(s,0,\frac{s}{2}\right)\)，\(E=\left(\frac{3s}{4},0,\frac{3s}{4}\right)\). 因而 \(\overrightarrow{DE}=\left(-\frac{s}{4},0,\frac{s}{4}\right)\)，\(\overrightarrow{AB_1}=(s,0,s)\)，两者内积为 \(0\)，所以 \(DE\perp AB_1\). 又 \(\overrightarrow{DC}=\left(-\frac{s}{2},h,-\frac{s}{2}\right)\)，有 \(\overrightarrow{DE}\cdot\overrightarrow{DC}=0\)，所以 \(DE\perp CD\). 直线 \(DE\) 分别在 \(E\)，\(D\) 处与 \(AB_1\)，\(CD\) 相交，故 \(DE\) 是两异面直线的公垂线.
\item 取 \(\overrightarrow{AB_1}=(s,0,s)\)，\(\overrightarrow{DC}=\left(-\frac{s}{2},h,-\frac{s}{2}\right)\). 由两直线夹角为 \(45^\circ\)，\(\frac{s^2}{s\sqrt{2}\sqrt{\frac{s^2}{2}+h^2}}=\cos45^\circ=\frac{1}{\sqrt{2}}\)，得 \(h^2=\frac{s^2}{2}\).

平面 \(A_1AC_1\) 和 \(A C_1B_1\) 的法向量可分别取 \(\bs n_1=\left(-\sqrt{2},1,0\right)\) 与 \(\bs n_2=\left(\sqrt{2},1,-\sqrt{2}\right)\). 设二面角为 \(\theta\)，则取其锐角，有 \(\cos\theta=\frac{|\bs n_1\cdot\bs n_2|}{|\bs n_1|\,|\bs n_2|}=\frac{1}{\sqrt{15}}\). 所以 \(\sin\theta=\sqrt{\frac{14}{15}}\)，从而 \(\tan\theta=\sqrt{14}\)，即二面角的大小为 \(\arctan\sqrt{14}\).
\end{enumerate}
\end{solution}

\begin{problem}
如图，由 \(M\) 到 \(N\) 的电路中有 \(4\) 个组件，分别标为 \(T_{1}\)，\(T_{2}\)，\(T_{3}\)，\(T_{4}\)，电流能通过 \(T_{1}\)，\(T_{2}\)，\(T_{3}\) 的概率都是 \(p\)，电流能通过 \(T_{4}\) 的概率是 \(0.9\). 电流能否通过各组件相互独立. 已知 \(T_{1}\)，\(T_{2}\)，\(T_{3}\) 中至少有一个能通过电流的概率为 \(0.999\).

\begin{enumerate}
\item 求 \(p\)；
\item 求电流能在 \(M\) 与 \(N\) 之间通过的概率.
\end{enumerate}

\bitmapfigure[max width=0.42\linewidth]{img/2010/syllabus_paper_2_liberal/q20_fig1.png}
\end{problem}

\begin{solution}
\begin{enumerate}
\item 设事件 \(A_i\) 表示电流能通过 \(T_i\). 由独立性，\(T_1\)，\(T_2\)，\(T_3\) 均不能通过的概率为 \((1-p)^3\). 因题设至少一个能通过的概率为 \(0.999\)，有 \((1-p)^3=1-0.999=0.001\). 因 \(0\leqslant p\leqslant1\)，得 \(1-p=0.1\)，所以 \(p=0.9\).
\item 从电路图看，电流通过 \(M\)，\(N\) 的事件为：\(T_4\) 能通过，或 \(T_4\) 不能通过且 \(T_3\) 能通过并且 \(T_1\)，\(T_2\) 中至少一个能通过. 这些情形互斥，因此所求概率为 \(0.9+0.1\cdot0.9\left[1-(1-0.9)^2\right]=0.9+0.1\cdot0.9\cdot0.99=0.9891\).
\end{enumerate}
\end{solution}

\begin{problem}
已知函数 \(f(x) = x^3 - 3ax^2 + 3x + 1\).

\begin{enumerate}
\item 设 \(a = 2\)，求 \(f(x)\) 的单调区间；
\item 设 \(f(x)\) 在区间 \((2,3)\) 中至少有一个极值点，求 \(a\) 的取值范围.
\end{enumerate}
\end{problem}

\begin{solution}
\begin{enumerate}
\item 当 \(a=2\) 时，\(f'(x)=3x^2-12x+3=3\left(x^2-4x+1\right)\). 方程 \(f'(x)=0\) 的根为 \(x=2-\sqrt{3}\) 和 \(x=2+\sqrt{3}\). 因二次项系数为正，\(f'(x)>0\) 在两根外侧成立，\(f'(x)<0\) 在两根之间成立. 所以增区间为 \(\left(-\infty,2-\sqrt{3}\right)\) 和 \(\left(2+\sqrt{3},+\infty\right)\)，减区间为 \(\left(2-\sqrt{3},2+\sqrt{3}\right)\).
\item \(f'(x)=3\left(x^2-2ax+1\right)\). 极值点必须是导数变号的零点. 对 \(x\in(2,3)\)，方程 \(f'(x)=0\) 等价于 \(a=\frac{x^2+1}{2x}=\frac{1}{2}\left(x+\frac{1}{x}\right)\). 函数 \(g(x)=\frac{1}{2}\left(x+\frac{1}{x}\right)\) 在 \((2,3)\) 上递增，因为 \(g'(x)=\frac{x^2-1}{2x^2}>0\). 因而存在 \(x\in(2,3)\) 使 \(f'(x)=0\) 的充要条件为 \(g(2)<a<g(3)\)，即 \(\frac{5}{4}<a<\frac{5}{3}\). 此时 \(x>2>1\)，另一个根为 \(1/x\)，两根不重合，导数确实变号，故该条件正好保证区间内至少有一个极值点.
\end{enumerate}
\end{solution}

\section{选考题：解答题}

\begin{problem}
已知斜率为 \(1\) 的直线 \(l\) 与双曲线 \(C: \frac{x^{2}}{a^{2}} - \frac{y^{2}}{b^{2}} = 1\)（\(a > 0\)，\(b > 0\)）相交于 \(B\)，\(D\) 两点，且 \(BD\) 的中点为 \(M(1,3)\).

\begin{enumerate}
\item 求 \(C\) 的离心率；
\item 设 \(C\) 的右顶点为 \(A\)，右焦点为 \(F\)，\(\left|DF\right| \cdot \left|BF\right| = 17\)，证明：过 \(A\)，\(B\)，\(D\) 三点的圆与 \(x\) 轴相切.
\end{enumerate}
\end{problem}

\begin{solution}
\begin{enumerate}
\item 直线 \(l\) 过 \(M(1,3)\) 且斜率为 \(1\)，所以 \(l\) 的方程为 \(y=x+2\). 设 \(B\)，\(D\) 的横坐标为 \(x_1\)，\(x_2\)，将直线方程代入双曲线方程，得 \((b^2-a^2)x^2-4a^2x-a^2b^2-4a^2=0\). 因 \(M\) 是 \(BD\) 的中点，\(x_1+x_2=2\). 由韦达定理，\(x_1+x_2=\frac{4a^2}{b^2-a^2}=2\)，所以 \(b^2=3a^2\). 双曲线的焦半距满足 \(c^2=a^2+b^2=4a^2\)，故离心率 \(e=\frac{c}{a}=2\).
\item 由 \(b^2=3a^2\)，双曲线方程等价于 \(3x^2-y^2=3a^2\). 设 \(B=(1+t,3+t)\)，\(D=(1-t,3-t)\)，其中 \(t>0\). 将 \(y=x+2\) 代入 \(3x^2-y^2=3a^2\)，得 \(2x^2-4x-4=3a^2\)，再令 \(x=1\pm t\)，得到 \(t^2=\frac{3(a^2+2)}{2}\).

右焦点为 \(F=(2a,0)\). 记 \(d_+^2=|BF|^2\)，\(d_-^2=|DF|^2\)，则 \(d_+^2=\left(1+t-2a\right)^2+\left(3+t\right)^2\)，\(d_-^2=\left(1-t-2a\right)^2+\left(3-t\right)^2\). 令 \(K=7a^2-4a+16\)，由上式和 \(2t^2=3a^2+6\) 得 \(d_+^2=K+4(2-a)t\)，\(d_-^2=K-4(2-a)t\). 因而 \(\left(|BF|\,|DF|\right)^2=K^2-16(2-a)^2t^2=25a^4+40a^3+96a^2+64a+64\).

由题设 \(|BF|\,|DF|=17\)，得到 \(25a^4+40a^3+96a^2+64a-225=0\)，即 \((a-1)\left(25a^3+65a^2+161a+225\right)=0\). 因 \(a>0\)，括号内三次式为正，故 \(a=1\). 此时右顶点 \(A=(1,0)\)，而 \(M=(1,3)\). 由于 \(BD\) 的斜率为 \(1\)，\(AM\) 为竖直线，故 \(AM\perp BD\). 又 \(M\) 是 \(BD\) 的中点，所以 \(A\) 在 \(BD\) 的垂直平分线上，即 \(AB=AD\). 因而过 \(A\)，\(B\)，\(D\) 的圆的圆心在直线 \(AM\) 上，圆心与 \(A\) 的连线垂直于 \(x\) 轴，所以该圆在 \(A\) 处与 \(x\) 轴相切.
\end{enumerate}
\end{solution}
